\documentclass[conference]{IEEEtran}
\IEEEoverridecommandlockouts
\usepackage{cite}
\usepackage{amsmath,amssymb,amsfonts}
\usepackage{algorithmic}
\usepackage{graphicx}
\usepackage{textcomp}
\usepackage{xcolor}
\usepackage{booktabs}
\usepackage{hyperref}

\def\BibTeX{{\rm B\kern-.05em{\sc i\kern-.025em b}\kern-.08em
    T\kern-.1667em\lower.7ex\hbox{E}\kern-.125emX}}

\begin{document}

\title{Task B1: Detecting Fraudulent Job Postings Using Multimodal Feature Fusion and Cost-Sensitive Classification\\
\large AI-101 Machine Learning Assignment --- Pool B: Mixed Modality (Numerical + Text)}

\author{\IEEEauthorblockN{Iman Fatima (Roll No: 70177713)}
\IEEEauthorblockA{\textit{Department of Intelligent Systems, Section C} \\
\textit{The University of Lahore}\\
Lahore, Pakistan \\
Email: iman.fatima@student.uol.edu.pk}
\and
\IEEEauthorblockN{Instructor: Dr. Muhammad Awais}
\IEEEauthorblockA{\textit{AI Architect} \\
\textit{The University of Lahore}\\
Lahore, Pakistan \\
Email: muhammad.awais@is.uol.edu.pk}
}

\maketitle

\begin{abstract}
Online employment fraud exposes job seekers to financial extortion and identity theft while degrading trust in recruitment platforms. Detecting fraudulent advertisements is challenging due to extreme class imbalance (4.84\% fraud rate, 19.65:1 ratio) and heterogeneous data comprising free-form job text and structured metadata. In this work, we develop an end-to-end multimodal classification pipeline on the EMSCAD benchmark (17,880 postings) utilizing scikit-learn \texttt{ColumnTransformer}. We evaluate balanced Logistic Regression, Linear SVM with modified Huber loss, and Random Forest across structured-only, text-only, and combined modalities. To address class imbalance, decision thresholds are calibrated on a dedicated validation split. On the held-out 15\% test set, the tuned Linear SVM achieves a fraud-class F1-score of 0.8760, Precision of 0.9464, and PR-AUC of 0.9110, outperforming single-modality models and establishing the efficacy of multimodal fusion.
\end{abstract}

\begin{IEEEkeywords}
Fraud detection, multimodal learning, class imbalance, TF-IDF, precision-recall AUC, cost-sensitive classification.
\end{IEEEkeywords}

\section{Introduction}
The proliferation of digital recruitment platforms has been accompanied by a surge in fraudulent job advertisements engineered for advance-fee scams, personal credential harvesting, and identity theft. Mitigating recruitment fraud requires automated, real-time screening systems that process diverse information streams including job titles, company profiles, unstructured position descriptions, and administrative metadata. In this study, we design and empirically evaluate a leak-free multimodal machine learning pipeline on the EMSCAD dataset (17,880 samples). Our primary finding is that joint multimodal feature fusion—integrating sublinear TF-IDF word unigrams/bigrams with metadata missingness indicators and categorical encodings—significantly outperforms unimodal baselines, boosting fraud-class F1 from 0.3119 (structured-only) and 0.8099 (text-only) to 0.8760 with an exceptional Precision-Recall AUC of 0.9110.

\section{Problem Statement}
We formulate fraudulent job-posting detection as an imbalanced binary classification problem. Each observation $i$ is represented by a composite feature vector $\mathbf{x}_i = [\mathbf{x}_i^{\text{text}}, \mathbf{x}_i^{\text{cat}}, \mathbf{x}_i^{\text{bin}}, \mathbf{x}_i^{\text{num}}]$, where $\mathbf{x}_i^{\text{text}}$ denotes unstructured natural language (title, company profile, description, requirements, benefits), $\mathbf{x}_i^{\text{cat}}$ represents categorical indicators (employment type, required experience, education), $\mathbf{x}_i^{\text{bin}}$ captures operational flags and missingness indicators, and $\mathbf{x}_i^{\text{num}}$ quantifies text length statistics. The objective is to learn a mapping $f: \mathbf{x} \to y \in \{0, 1\}$, where $y=1$ denotes a fraudulent advertisement and $y=0$ denotes a legitimate posting. Because the positive class comprises only 4.84\% of instances, accuracy is clinically misleading (a naive majority-class rule achieves 95.16\% accuracy with 0\% recall). Performance is governed by fraud-class Precision, Recall, F1-score, and Precision-Recall Area Under the Curve (PR-AUC).

\section{Dataset Used}
We employ the Employment Scam Aegean Dataset (EMSCAD) curated by the University of the Aegean, hosted on Kaggle~\cite{emscad}. The benchmark contains 17,880 genuine and fake job descriptions collected between 2012 and 2014, spanning 18 attributes. The target attribute \texttt{fraudulent} divides into 17,014 legitimate (95.16\%) and 866 fraudulent (4.84\%) postings. Missing values are prevalent: 84.0\% of records lack salary information, 40.5\% lack department designations, 18.5\% lack company profiles, and 15.1\% lack explicit requirements. Rather than imputing or discarding missing entries, we engineer explicit missingness indicator flags, transforming missing data into predictive signals.

\section{Dataset Statistics}
Table~\ref{tab:stats} summarizes the class distribution and text characteristics across the dataset. Fraudulent postings exhibit significantly shorter narratives: the median word count for fraudulent postings is 224.5 words, compared to 355.0 words for authentic postings.

\begin{table}[htbp]
\caption{Dataset Modality & Class Breakdown (17,880 Samples)}
\label{tab:stats}
\centering
\small
\begin{tabular}{lrrr}
\toprule
\textbf{Class / Partition} & \textbf{Count} & \textbf{Percentage} & \textbf{Median Words} \\
\midrule
Legitimate (Real = 0) & 17,014 & 95.16\% & 355.0 \\
Fraudulent (Fake = 1) & 866 & 4.84\% & 224.5 \\
\midrule
Training Set (70\%) & 12,516 & 70.00\% & 606 Frauds \\
Validation Set (15\%) & 2,682 & 15.00\% & 130 Frauds \\
Held-out Test Set (15\%) & 2,682 & 15.00\% & 130 Frauds \\
\bottomrule
\end{tabular}
\end{table}

Exploratory analysis demonstrates that administrative metadata correlates strongly with fraud risk: postings lacking a verified corporate logo (\texttt{has\_company\_logo=0}) carry a 15.93\% fraud rate versus 1.99\% when present (an 8.0$\times$ disparity). Similarly, missing company profiles carry a 17.74\% fraud rate versus 1.91\% when present (a 9.3$\times$ increase).

\section{Methodology}

\subsection{Learning Setting}
We deploy a unified scikit-learn \texttt{ColumnTransformer} architecture to ingest heterogeneous modalities without data leakage:
\begin{itemize}
    \item \textbf{Text Pipeline:} Unifies five text fields, applying lowercasing, regex cleaning, and \texttt{TfidfVectorizer} with sublinear term-frequency scaling, unigram/bigram tokens ($n \in \{1, 2\}$), English stop-word removal, and vocabulary capped at 10,000 features.
    \item \textbf{Categorical Pipeline:} Enforces one-hot encoding on \texttt{employment\_type}, \texttt{required\_experience}, and \texttt{required\_education} with unseen-category handling.
    \item \textbf{Binary \& Missingness Pipeline:} Passes binary operational flags (\texttt{telecommuting}, \texttt{has\_company\_logo}, \texttt{has\_questions}) and engineered missingness indicators (\texttt{missing\_profile}, \texttt{missing\_salary}, etc.).
    \item \textbf{Numeric Pipeline:} Applies \texttt{StandardScaler} to character and word counts.
\end{itemize}

We benchmark three model families equipped with balanced class weighting ($\frac{N}{2 \cdot N_k}$): (1) Logistic Regression baseline; (2) Linear SVM optimized via SGD with modified Huber loss for calibrated posterior probability estimates; and (3) Random Forest Classifier ($N_{\text{est}}=150$, depth$=25$).

\subsection{Data Split Approach}
The 17,880 observations were divided using a stratified 70/15/15 ratio into 12,516 training, 2,682 validation, and 2,682 test instances with fixed seed 42. Preprocessing statistics were strictly fitted on the training split. The validation partition was reserved exclusively for decision threshold tuning. The test split remained completely untouched until final evaluation.

\subsection{Parameter Settings}
Decision thresholds were optimized on the validation set by sweeping $\tau \in [0.10, 0.90]$ to maximize fraud-class F1. Table~\ref{tab:hyperparams} documents model parameters and selected thresholds.

\begin{table}[htbp]
\caption{Model Hyperparameters and Decision Thresholds}
\label{tab:hyperparams}
\centering
\small
\begin{tabular}{lll}
\toprule
\textbf{Model} & \textbf{Hyperparameters} & \textbf{Optimal $\tau$} \\
\midrule
Logistic Regression & $C=1.0$, balanced class weight, L2 penalty & $\tau = 0.73$ \\
Linear SVM (Huber) & Modified Huber loss, $\alpha = 10^{-4}$, L2 penalty & $\tau = 0.71$ \\
Random Forest & $N_{\text{est}}=150$, Max Depth $=25$, balanced & $\tau = 0.68$ \\
\bottomrule
\end{tabular}
\end{table}

\section{Results and Analysis}

\subsection{Results}
Table~\ref{tab:results} summarizes the test-set performance across candidate models evaluated at their optimal decision thresholds.

\begin{table}[htbp]
\caption{Final Test-Set Evaluation Metrics ($\tau$ Optimized)}
\label{tab:results}
\centering
\small
\begin{tabular}{lrrrrr}
\toprule
\textbf{Model} & \textbf{Prec.} & \textbf{Recall} & \textbf{F1} & \textbf{PR-AUC} & \textbf{ROC-AUC} \\
\midrule
M1: Logistic Regression & 0.8271 & \textbf{0.8462} & 0.8365 & 0.9067 & \textbf{0.9913} \\
M2: Linear SVM (Huber)  & \textbf{0.9464} & 0.8154 & \textbf{0.8760} & \textbf{0.9110} & 0.9799 \\
M3: Random Forest       & 0.8736 & 0.5846 & 0.7005 & 0.8256 & 0.9832 \\
\bottomrule
\end{tabular}
\end{table}

\subsection{Analysis}
\textbf{1. Why Accuracy is a Misleading Metric:}
Due to extreme 19.65:1 class imbalance, an uninformative dummy predictor assigning every posting to "Legitimate" achieves an apparent accuracy of 95.16\%, yet yields 0.00\% recall and 0.00 F1 on scams. In high-stakes fraud filtering, false negatives allow cybercriminals to exploit vulnerable applicants, while false positives needlessly suppress legitimate recruitment ads. Evaluating models using Precision-Recall curves (PR-AUC) and positive-class F1-score isolates performance on rare positive events without inflation from the dominant true negative class.

\textbf{2. Multimodal Ablation Study:}
Table~\ref{tab:ablation} presents the mandatory ablation comparing structured metadata only, text features only, and joint multimodal fusion on identical test data using Logistic Regression.

\begin{table}[htbp]
\caption{Feature Modality Ablation Study on Test Set}
\label{tab:ablation}
\centering
\small
\begin{tabular}{lrrrrr}
\toprule
\textbf{Modality} & \textbf{Prec.} & \textbf{Recall} & \textbf{F1} & \textbf{PR-AUC} & \textbf{ROC-AUC} \\
\midrule
A1: Structured-Only & 0.2137 & 0.5769 & 0.3119 & 0.2244 & 0.8461 \\
A2: Text-Only       & \textbf{0.8750} & 0.7538 & 0.8099 & 0.8943 & 0.9873 \\
A3: Combined Fusion & 0.8271 & \textbf{0.8462} & \textbf{0.8365} & \textbf{0.9067} & \textbf{0.9913} \\
\bottomrule
\end{tabular}
\end{table}

The ablation proves that joint multimodal fusion substantially outperforms either modality alone. Structured features alone lack semantic depth (F1 = 0.3119). Text features provide strong semantic discrimination (F1 = 0.8099). Fusing text with structured flags increases fraud recall from 75.38\% to 84.62\% and raises PR-AUC from 0.8943 to 0.9067 (+0.0124), demonstrating that missingness indicators and operational flags provide complementary orthogonal signals.

\textbf{3. Discriminative Lexical and Metadata Signals:}
Inspection of model feature weights reveals that terms exhibiting the largest positive weights (indicative of fraud) include \textit{aptitude staffing} (+3.742), \textit{earn} (+3.078), \textit{money} (+2.985), \textit{data entry} (+2.862), \textit{phone} (+2.750), and \textit{signing bonus} (+2.479). These tokens reflect classic work-from-home financial transfer and data-entry phishing scams. Conversely, authentic postings are governed by high word count (-2.367), \textit{team} (-1.848), \textit{client} (-1.812), \textit{software} (-1.575), and \textit{join team} (-1.436).

Analysis of false positives reveals legitimate early-stage startups that omit corporate logos and profiles, triggering high missingness penalties. Adding company verification checks mitigates these edge cases.

\section{Conclusion}
This study established an effective multimodal machine learning architecture for detecting fraudulent job postings under extreme class imbalance. Linear SVM with modified Huber loss combined with calibrated decision thresholding achieved a fraud-class F1 of 0.8760 and a PR-AUC of 0.9110 on held-out test data. Multimodal fusion of text and structured metadata yielded an 11-point recall boost over text-only models. Future work will investigate fine-tuning contextual transformer encoders (DeBERTa-v3) with focal loss to further separate subtle fraudulent linguistic patterns.

\section*{AI Use Disclosure}
I used Google DeepMind Antigravity AI assistant for structuring the multimodal feature extraction pipeline, constructing scikit-learn ColumnTransformer abstractions, and formatting the IEEE conference report. All dataset preprocessing, threshold calibration, ablation experiments, metric audits, and written analyses were produced, verified, and checked by me, and I take full responsibility for this work.

\begin{thebibliography}{00}
\bibitem{emscad} S. Vidros, C. Kolias, G. Kambourakis, and L. Akoglu, ``Automatic Detection of Online Recruitment Frauds: Characteristics, Methods, and a Public Dataset,'' \textit{Future Internet}, vol. 9, no. 1, p. 6, 2017.
\bibitem{kaggle} S. Bansal, ``Real or Fake: Fake JobPosting Prediction Dataset,'' Kaggle, 2020.
\bibitem{scikit-learn} F. Pedregosa \textit{et al.}, ``Scikit-learn: Machine learning in Python,'' \textit{Journal of Machine Learning Research}, vol. 12, pp. 2825--2830, 2011.
\end{thebibliography}

\end{document}
